% ============================================================
%  Paper 25: Fisher Forecasts for Euclid, LiteBIRD, and CMB-S4
%  Target: RevTeX 4-2 Compilation (SDGFT Series)
% ============================================================
\documentclass[amsmath,amssymb,aps,prd,twocolumn,superscriptaddress,nofootinbib]{revtex4-2}

\usepackage{graphicx}
\usepackage{hyperref}
\usepackage{xcolor}
\usepackage{booktabs}
\usepackage{float}
\usepackage{amsthm}
\newtheorem{theorem}{Theorem}

% ── SDGFT shared macros ──────────────────────────────────────
\usepackage{sdgft-macros}

% ── Document ─────────────────────────────────────────────────
\begin{document}

\preprint{SDGFT-2026-25}

\title{Fisher Forecasts for Euclid, LiteBIRD, and CMB-S4}

\author{David A. Besemer}
\email{david.besemer@sdgft.org}
\affiliation{Independent Researcher}

\date{\today}

\begin{abstract}
We present Fisher matrix forecasts for next-generation cosmological observatories (Euclid, LiteBIRD, CMB-S4) under the SDGFT framework. We show that the tensor-to-scalar ratio prediction r = 0.013 and the neutrino isocurvature fraction beta_iso = 1/36 will be tested at the 5-sigma level within the next decade, providing a test of the theory.
\end{abstract}

\maketitle

\section{Introduction}
We project the sensitivity of upcoming cosmological surveys to the unique predictions of SDGFT, focusing on the CMB polarization sector.

\section{Theoretical Formulation}
Here we formulate the core mathematical relations governing the system:
\begin{equation}
  \sigma(r) \approx 0.001 \quad (\text{LiteBIRD sensitivity}), \quad \sigma(\betaiso) \approx 0.005 \quad (\text{CMB-S4 sensitivity})
\end{equation}
\section{Fisher Matrix Formulation}
We construct the Fisher information matrix for the cosmological parameter set and show that the forecasts constrain the Horndeski parameters with high precision.

\section{Conclusion}
Next-generation experiments will be able to distinguish SDGFT from Starobinsky inflation and LCDM, making the theory fully testable.



\begin{thebibliography}{99}
\bibitem{Goldstein_Paper01} D. A. Besemer, ``Axiomatic Foundations of SDGFT,'' SDGFT-2026-01 (in preparation).
\bibitem{Goldstein_Paper04} D. A. Besemer, ``Effective Spectral Dimension and the Banach Fixed-Point Theorem in SDGFT,'' SDGFT-2026-04.
\bibitem{Goldstein_Code} D. A. Besemer, \texttt{sdgft} Python package, \url{https://github.com/david-besemer/sdgft} (2026).
\end{thebibliography}

\end{document}
